% Converted from the authoritative recovery source, 2026-11-14.
\subsection{Why both additions matter}\label{influence-recovery:why-both-additions-matter}

For magnetization \(M=n^{-1}\sum_iX_i\) in zero-field Curie--Weiss,
exchangeability and sign symmetry imply \[
 I_i(M)=n^{-2},\qquad
 \operatorname{Var}(\mathbb E[M\mid X_i])
   =\operatorname{Var}(M)^2.
\] The reveal term captures ordered-phase magnetization with a single
query. But merely adding singleton reveal gain to flip influence does
not work for every exchangeable law. Let \[
 \mu=\tfrac12U_n+\tfrac14\delta_{(+1)^n}
                       +\tfrac14\delta_{(-1)^n},\qquad f=M^2.
\] This law has full support. Its \(f\)-variance tends to \(1/4\), while
\(I_i(f)\le4/n^2\). Global sign symmetry makes \(\mathbb E[f\mid X_i]\)
constant, so every singleton reveal gain is zero. A short prefix can
identify the density sector even though a single observation cannot;
derivative influences handle variance remaining inside a sector. This
explains the bounded lookahead block.

The theorem is derived from exchangeability and UAA, not asserted for
all dependent Gibbs laws. A law with a genuine quantum-versus-classical
query separation cannot satisfy a comparable hereditary polynomial
progress theorem for all relevant bounded polynomials.

\subsubsection{A full-support Gibbs example with bounded one-spin energy changes}
\label{influence-recovery:smooth-singleton-obstruction}

The independent audit strengthens the preceding singleton obstruction.
Fix \(0<r<1\), and draw a latent mean \(\Theta\) equal to zero with
probability \(1/2\), and to each of \(+r,-r\) with probability \(1/4\).
Conditional on \(\Theta\), draw independent spins with mean \(\Theta\).
Its density relative to the uniform law is
\[
 w_n(x)=\frac12+\frac14\prod_i(1+rx_i)
                    +\frac14\prod_i(1-rx_i).
\]
This is a full-support exchangeable law and is the Gibbs law of
\(H_n=-\log w_n\) at inverse temperature one. Its energy range is \(O(n)\).
A one-spin change alters the energy by at most
\(\log((1+r)/(1-r))\): the ratio of every positive summand before and
after a change lies in the corresponding fixed interval, as does the
ratio of their sums.

Let \(M=n^{-1}\sum_i x_i\) and \(f=M^2\in[0,1]\), of degree two.
Conditional laws of large numbers give
\[
 \operatorname{Var}(f)\longrightarrow r^4/4>0.
\]
Meanwhile
\[
 D_i f=\frac{2\sum_{j\ne i}x_j}{n^2},\qquad I_i(f)\le4/n^2.
\]
The law is invariant under global spin reversal and \(f\) is even, so
its two conditional expectations given \(X_i=+1\) and \(X_i=-1\) agree:
\[
 \operatorname{Var}(\mathbb E[f\mid X_i])=0
 \quad\text{for every }i.
\]
Every temperature of this Hamiltonian remains exchangeable and hence
belongs to the all-temperature approximation class proved above.
Adding the two singleton quantities cannot provide a general influence
repair. A short transcript can reveal higher-order phase information;
recoverable-source influence gives another possible repair.
